\section{Equazioni di Hamilton}

Le equazioni trattate in questa sezione vengono dette \emph{equazioni di Hamilton} o \emph{equazioni canoniche della meccanica} e il loro sviluppo è datato al 1834.

\subsection{Trasformata di Legendre}

Consideriamo il calcolo del differenziale di una funzione di due variabili $f(x,y)$:
\begin{equation*}
    \dd{f} = \pdv{f}{x} \dd{x} + \pdv{f}{y} \dd{y} = u \dd{x} + v \dd{y}.
\end{equation*}
Talvolta può risultare utile scrivere il differenziale di $f$ in termini di $u$ e $v$ invece che di $x$ e $y$. Attraverso la regola di Leibnitz si ha
\begin{equation*}
    \dd{f} = \dd{(ux)} - x \dd{u} + \dd{(vy)} - y \dd{v} = \dd{(ux)} - x \dd{u} + v \dd{y},
\end{equation*}
e portando al primo membro si ha
\begin{equation*}
    \dd{(f - ux)} = - x \dd{u} + v \dd{y}.
\end{equation*}
Ponendo $g(u,y) \defeq f(x,y) - ux$ si ha che
\begin{equation*}
    \dd{g} = - x \dd{u} + v \dd{y},
\end{equation*}
avendo invertito la dipendenza di $f$ da $x$ a $u$, ottenendo la \emph{trasformata di Legendre} di $f$.

\begin{definition}[Trasformata di Legendre]
Sia $f(x,y)$ una funzione di due variabili. La \emph{trasformata di Legendre} di $f$ è la funzione $g(u,y) = f(x,y) - ux$ ottenuta invertendo la dipendenza di $f$ da $x$ a $u$.
    \begin{equation*}
        u = \pdv{f}{x} \iff x = - \pdv{g}{u}.
    \end{equation*}
\end{definition}


\subsection{Derivazione delle equazioni di Hamilton}
Partiamo dalle equazioni di Lagrange
\begin{equation*}
    \dv{t} \pdv{L}{\dot{q}_i} - \pdv{L}{q_i} = 0.
\end{equation*}
Le equazioni possono essere riscritte in termini dei momenti coniugati
\begin{equation*}
    p_i = \pdv{L}{\dot{q}_i}, \qquad \dot{p_i} = \pdv{L}{q_i}, \tag{$*$}
\end{equation*}
ottenendo $2n$ equazioni differenziali del primo ordine invece di $n$ equazioni differenziali del secondo ordine.

Applichiamo la trasformata di Legendre alla lagrangiana
\begin{align*}
    \dd{L} &= \sum_{j=1}^n \left( \pdv{L}{\dot{q}_j} \dd{\dot{q}_j} - \pdv{L}{q_j} \dd{q_j} \right) \\
    &= \sum_{j=1}^n \dot{p_j} \dd{q_j} +  \sum_{j=1}^n p_j \dd{\dot{q}_j} + \pdv{L}{t} \dd{t} \\
    &= \sum_{j=1}^n \dot{p_j} \dd{q_j} +  \dd {\left(\sum_{j=1}^n p_j \dot{q}_j\right)} - \sum_{j=1}^n \dot{q_j} \dd{p_j} + \pdv{L}{t} \dd{t}.
\end{align*}
Portando al primo membro si ha
\begin{equation*}
    \dd \left(\sum_{j=1}^n p_j\dot{q}_j - L\right) = \sum_{j=1}^n \dot{p_j} \dd{q_j} - \sum_{j=1}^n \dot{q_j} \dd{p_j} + \pdv{L}{t} \dd{t}.
\end{equation*}
Definiamo quindi la funzione
\begin{equation*}
    H(\vq, \vp, t) \defeq \sum_{j=1}^n p_j\dot{q}_j - L(\vq, \dot{\vq}, t),
\end{equation*}
ottenendo
\begin{equation*}
    -\dd{H} = \sum_{j=1}^n \dot{p_j} \dd{q_j} - \sum_{j=1}^n \dot{q_j} \dd{p_j} + \pdv{L}{t} \dd{t}.
\end{equation*}
Il differenziale di $H$ è
\begin{equation*}
    \dd{H} = \sum_{j=1}^n \pdv{H}{q_j} \dd{q_j} + \sum_{j=1}^n \pdv{H}{p_j} \dd{p_j} + \pdv{H}{t} \dd{t},
\end{equation*}
e confrontando i due differenziali si ha
\begin{equation*}
 \dot{q}_j = \pdv{H}{p_j}, \qquad \dot{p}_j = -\pdv{H}{q_j}, \qquad \pdv{H}{t} = -\pdv{L}{t}.
\end{equation*}

La funzione $H$ è la trasformata di Legendre di $L$ e viene detta \emph{hamiltoniana} del sistema.

Risolvendo le equazioni $(*)$ per $\dot{q}_j$ si ha che $H$ è una funzione di $\vq$, $\vp$ e $t$, e non dipende più da $\dot{\vq}$, ottenendo così una trasformazione che porta da un sistema di coordinate $(\vq, \dot{\vq})$ ad un sistema di coordinate $(\vq, \vp)$:
\begin{equation*}
    (\vq, \dot{\vq}) \leftrightarrow (\vq, \vp),
\end{equation*}
passando dallo spazio $n$-dimensionale delle configurazioni $\mathcal{C}$ allo spazio $2n$-dimensionale detto \emph{spazio delle fasi} di Gibbs $\Phi_{2n}$.

Le $2n$ equazioni differenziali ordinarie del primo ordine in $2n$ incognite
\begin{equation*}
    \begin{cases}
        \dot{q}_j = \pdv{H}{p_j} \\
        \dot{p}_j = -\pdv{H}{q_j}
    \end{cases}
\end{equation*}
vengono dette \emph{equazioni di Hamilton} o \emph{equazioni canoniche della meccanica}.

Nel caso meccanico, in un sistema naturale a vincoli scleronomi, la lagrangiana è data da $L = T - V$, e quindi
\begin{equation*}
    p_i = \pdv{L}{\dot{q}_i} = \pdv{T}{\dot{q}_i} = \pdv{T_2}{\dot{q}_j} = \sum_{j=1}^n a_{ij} \dot{q}_j.
\end{equation*}
Otteniamo quindi un sistema lineare di equazioni
\begin{equation*}
    \vp = A \dot{\vq}.
\end{equation*}
Sostituendo nella definizione di $H$ si ha
\begin{equation*}
    H = \sum_{j=1}^n p_j \dot{q}_j - L = \sum_{i,j=1}^n a_{ij} \dot{q}_i \dot{q}_j - (T - V) = T + V,
\end{equation*}
ovvero l'hamiltoniana coincide con l'energia totale del sistema.

\begin{example}[Oscillatore armonico]
    % Guida orizzontale, molla con coefficiente k nell'origine punto P di massa M
    Consideriamo un oscillatore armonico, ovvero una massa $m$ che si muove su una guida orizzontale e che è collegata ad una molla con coefficiente elastico $k$ ancorata all'origine. La coordinata generalizzata è data da $q = x$, e la lagrangiana è data da
    \begin{equation*}
        L = T - V = \frac{1}{2} m \dot{q}^2 - \frac{1}{2} k q^2.
    \end{equation*}

    Il momento coniugato è dato da
    \begin{equation*}
        p = \pdv{L}{\dot{q}} = m \dot{q},
    \end{equation*}
    e l'hamiltoniana è data da
    \begin{equation*}
        H = \left[p \dot{q} - L\right]_{\dot{q} = \frac{1}{m} p} = \frac{p^2}{2m} + \frac{1}{2} k q^2 = \frac{1}{2m}\left(p^2 + mk q^2\right) = \frac{1}{2m}\left(p^2 + m^2\omega^2 q^2\right).
    \end{equation*}
    Le equazioni di Hamilton sono date da
    \begin{equation*}
        \begin{cases}
            \dot{q} = \pdv{H}{p} = \frac{p}{m} \\
            \dot{p} = -\pdv{H}{q} = m\omega^2 q
        \end{cases}
    \end{equation*}
\end{example}

Va notato che la funzione $H$ presenta una forma estremamente simile alla funzione $h$ di Jacobi:
\begin{equation*}
    H(\vq, \vp, t) = \sum_{j=1}^n p_j \dot{q}_j - L(\vq, \dot{\vq}, t), \qquad h(\vq, \dot{\vq}, t) = \sum_{j=1}^n \pdv{L}{\dot{q}_j} \dot{q}_j - L(\vq, \dot{\vq}, t).
\end{equation*}
La differenza però è sostanziale: $H$ è una funzione di $\vq$, $\vp$ e $t$, mentre $h$ è una funzione di $\vq$, $\dot{\vq}$ e $t$.

\subsection{Variabili cicliche in un sistema hamiltoniano}

\begin{definition}[Coodinata ciclica in un sistema hamiltoniano]\note{Lezione 23/04/2026}
    Una coordinata $q_i$ dicesi \emph{ciclica} per il sistema di hamiltoniana $H(\vq, \vp, t)$ se 
    \begin{equation*}
        \pdv{H}{q_i} = 0.
    \end{equation*}
\end{definition}
Derivando la relazione di definizione di $H$ rispetto a $q_i$ si ha
\begin{equation*}
    \pdv{H}{q_i} = - \pdv{L}{q_i},
\end{equation*}
per cui se $q_i$ è ciclica per $H$ allora è ciclica anche per $L$.

Ipotizziamo ora che in un dato sistema hamiltoniano tutte le coordinate $q_i$ siano cicliche e che $H$ non dipenda esplicitamente dal tempo:
\begin{equation*}
    \pdv{H}{q_i} = 0, \qquad i = 1, \dots, n.
\end{equation*}
In questo caso l'hamiltoniana sarebbe una funzione delle sole $p_i$ e le equazioni di Hamilton sarebbero date da
\begin{equation*}
    \begin{cases}
        \dot{q}_i = \pdv{H}{p_i} \\
        \dot{p}_i = 0
    \end{cases}
\end{equation*}
per cui i momenti coniugati $p_i$ sarebbero costanti 
\begin{equation*}
    p_i = \text{costante} = \alpha_i, \qquad i = 1, \dots, n.
\end{equation*}
Definiamo la quantità
\begin{equation*}
    \omega_i = \left. \pdv{H}{p_i} \right|_{p_i = \alpha_i},
\end{equation*}
e integrando le equazioni di Hamilton si ha
\begin{align*}
    \int \dot{q}_i \dd{t} &= \int \omega_i \dd{t} \\
    \implies q_i &= \omega_i t + \beta_i, \qquad i = 1, \dots, n.
\end{align*}
Si ha quindi che se è possibile trovare un sistema di coordinate in cui tutte le coordinate siano cicliche allora il sistema è risolvibile per \emph{quadrature}, ovvero è possibile risolvere le equazioni di Hamilton attraverso un numero finito di integrazioni.

Considerando un moto del sistema $\vq(t)$, $\vp(t)$ e mantenendo l'ipotesi che $\pdv{H}{t} = 0$, calcoliamo
\begin{align*}
    \dv{H(\vq(t), \vp(t))}{t} &= \sum_{j=1}^n \left( \pdv{H}{q_j} \dot{q}_j + \pdv{H}{p_j} \dot{p}_j \right) + \pdv{H}{t}
    \\ &= \sum_{j=1}^n \left( \pdv{H}{q_j} \pdv{H}{p_j} - \pdv{H}{p_j} \pdv{H}{q_j} \right) \\
    &= \pdv{H}{t} = 0.
\end{align*}

\subsection{Principio di Hamilton modificato}
Partiamo dal principio di Hamilton e sostituiamo l'inversa della trasformazione di Legendre della lagrangiana:
\begin{equation*}
    \var \int_{t_0}^{t_1} L(\vq, \dot{\vq}, t) \dd{t} = \var \int_{t_0}^{t_1} \left( \sum_{j=1}^n p_j \dot{q}_j - H(\vq, \vp, t) \right) \dd{t} = 0.
\end{equation*}
Applicando l'operatore differenziale si ha
\begin{equation*}
\var \int_{t_0}^{t_1} L \dd{t} = \int_{t_0}^{t_1} \sum_{k=1}^n \left( \pdv{L}{q_k} - \dv{t}\pdv{L}{\dot{q}_k} \right) \var{q_k} \dd{t} + \left[\sum_{k=1}^n \pdv{L}{\dot{q}_k} \var{q}_k \right]_{t_0}^{t_1}.
\end{equation*}
Considerando ora come lagrangiana la funzione
\begin{equation*}
    L^*(\vq, \vp, \dot{\vq}, \dot{\vp}, t) = \sum_{j=1}^n p_j \dot{q}_j - H(\vq, \vp, t),
\end{equation*}
otteniamo
\begin{multline*}
    \var \int_{t_0}^{t_1} L^* \dd{t}
    = \int_{t_0}^{t_1} \sum_{k=1}^n \left( \pdv{L^*}{q_k} - \dv{t}\pdv{L^*}{\dot{q}_k} \right) \var{q_k} \dd{t}  + \int_{t_0}^{t_1} \sum_{j=1}^n \left( \pdv{L^*}{p_j} - \dv{t}\pdv{L^*}{\dot{p}_j} \right) \var{p_j} \dd{t} + \\
    + \left[\sum_{j=1}^n \pdv{L^*}{\dot{q}_j} \var{q_j}+ \sum_{j=1}^n \pdv{L^*}{\dot{p}_j} \var{p_j} \right]_{t_0}^{t_1} = 0.
\end{multline*}
Notiamo che le sommatorie negli ultimi due termini si annullano, in quanto $\var{q_j}(t_0) = \var{q_j}(t_1) = 0$ e $\pdv{L^*}{\dot{p}_j} = 0$ per ogni $j$.
Quindi, le equazioni di Eulero-Lagrange per $L^*$ sono
\begin{align*}
    \pdv{L^*}{q_j} - \dv{t}\pdv{L^*}{\dot{q}_j} &= - \pdv{H}{q_j} - \dot{p}_j = 0 \iff \pdv{H}{q_j} = - \dot{p}_j \\
    \pdv{L^*}{p_j} - \dv{t}\pdv{L^*}{\dot{p}_j} &= \dot{q}_j - \pdv{H}{p_j} = 0 \iff \pdv{H}{p_j} = \dot{q}_j
\end{align*}
ottenendo così le equazioni di Hamilton.

\subsection{Parentesi di Poisson}
\begin{definition}[Parentesi di Poisson]
    Siano $f(\vq, \vp, t)$ e $g(\vq, \vp, t)$ due funzioni di classe $C^1$ nelle variabili canoniche\note{Le variabili canoniche sono le coordinate generalizzate $q_i$ e i momenti coniugati $p_i$}. Si definisce \emph{parentesi di Poisson} di $f$ e $g$
    \begin{equation*}
        \{f, g\} \defeq \sum_{j=1}^n \left( \pdv{f}{q_j} \pdv{g}{p_j} - \pdv{f}{p_j} \pdv{g}{q_j} \right).
    \end{equation*}
\end{definition}

Analizziamo ora le proprietà delle parentesi di Poisson. Siano $f(\vq, \vp, t)$, $g(\vq, \vp, t)$ e $h(\vq, \vp, t)$ tre funzioni di classe $C^1$ nelle variabili canoniche, allora
\begin{enumerate}[label=($\roman*$)]
    \item $\{f, f\} = 0$;
    \item $\{f, g\} = - \{g, f\}$ (antisimmetria);
    \item $\{c, f\} = 0$, per ogni $c\in\R$;
    \item $\{\alpha f + \beta g, h\} = \alpha \{f, h\} + \beta \{g, h\}$, per ogni $\alpha, \beta\in \R$  (linearità);
    \item $\pdv{t} \{f, g\} = \{\dot{f}, g\} + \{f, \dot{g}\}$;
    \item $\{f, \{g, h\}\} + \{g, \{h, f\}\} + \{h, \{f, g\}\} = 0$ (identità di Jacobi--Poisson). \note{permutando ciclicamente le funzioni}
\end{enumerate}

\vspace{1em}
È interessante notare che la proprietà ($v$) è valida considerando qualsiasi operatore differenziale $\mathcal{D}$ al posto della derivata temporale $\pdv{t}$:
\begin{equation*}
    \mathcal{D} \{f, g\} = \{\mathcal{D} f, g\} + \{f, \mathcal{D} g\}.
\end{equation*}

\begin{definition}[Integrale primo per le equazioni di Hamilton]
    Una funzione $f(\vq, \vp, t)$ è un \emph{integrale primo} per le equazioni di Hamilton se $f$ si conserva costante sulle soluzioni di tali equazioni.
\end{definition}
\begin{remark}
    $f$ è un integrale primo se e soltanto se 
    \begin{equation*}
        \pdv{f}{t} + \{f, H\} = 0.
    \end{equation*}
    Infatti data una soluzione $\vq(t)$, $\vp(t)$ delle equazioni di Hamilton, si ha che
    \begin{align*}
        \dv{f(\vq(t), \vp(t), t)}{t} &= \pdv{f}{t} + \sum_{j=1}^n \left( \pdv{f}{q_j} \dot{q}_j + \pdv{f}{p_j} \dot{p}_j \right) \\
        &= \pdv{f}{t} + \sum_{j=1}^n \left(\pdv{f}{q_j} \pdv{H}{p_j} - \pdv{f}{p_j} \pdv{H}{q_j} \right) \\
        &= \pdv{f}{t} + \{f, H\}.
    \end{align*}
\end{remark}

\begin{theorem}[di Jacobi--Poisson]
    Se $f$ e $g$ sono integrali primi delle equazioni di Hamilton, allora anche $\{f, g\}$ è un integrale primo.
\end{theorem}
\begin{proof}
    Vogliamo dimostrare che 
    \begin{equation*}
        \{f,g\} \text{ è un integrale primo} \iff \pdv{t} \{f, g\} + \{\{f, g\}, H\} = 0.
    \end{equation*}
    Calcoliamo
    \begin{align*}
        \pdv{t} \{f, g\} + \{\{f, g\}, H\} &= \{\dot{f}, g\} + \{f, \dot{g}\} + \{\{f, g\}, H\} \\
        &= -\{\{f,H\}, g\} - \{f, \{g, H\}\} - \{H, \{f, g\}\} \\
        &= -\{g, \{H, f\}\} - \{f, \{g, H\}\} - \{H, \{f, g\}\} \\
        &= 0,
    \end{align*}
    avendo utilizzato l'identità di Jacobi--Poisson.
\end{proof}

Le seguenti relazioni vengono dette \emph{parentesi di Poisson fondamentali}:
\begin{itemize}
    \item $\{q_i, q_j\} = 0$;
    \item $\{p_i, p_j\} = 0$;
    \item $\{q_i, p_j\} = \delta_{ij}$,
\end{itemize}
dove $\delta_{ij}$ è la delta di Kronecker.